\NativeTitle{105}{期末综合素养评价（四）}{}{90}{5}{一 & 二 & 三 & 四 & 五}{ &  &  &  & }
\NativeSlot{277}{41}{%
\ExamGroup{一、填空。{\normalfont\kaishu\fontsize{10.8bp}{14bp}\selectfont （每空1分，共21分）}}

% p105_s01

\ExamQuestion{1.}{5}%
\(9.08\)读作\mbox{（\hspace{7em}）}，它的计数单位是\mbox{（\hspace{4.5em}）}，它有\mbox{（\hspace{1.8em}）}个这样的计数单位，再添上\mbox{（\hspace{1.8em}）}个这样的计数单位就是最小的两位数。\par
}
\NativeSlot{351}{36}{%
% p105_s02

\ExamQuestion{2.}{5}%
\(6\)千克\(50\)克＝\mbox{（\hspace{1.8em}）}克；\allowbreak{}\(15.6\)元＝\mbox{（\hspace{1.8em}）}元\mbox{（\hspace{1.8em}）}角；\allowbreak{}\(900\)克＝\mbox{（\hspace{1.8em}）}千克；\allowbreak{}\(3.07\)千米＝\mbox{（\hspace{1.8em}）}千米\mbox{（\hspace{1.8em}）}米。\par
}
\NativeSlot{401}{36}{%
% p105_s03

\ExamQuestion{3.}{5}%
用\(1,0,2,7\)和小数点组成一个三位小数，这个小数最大是\mbox{（\hspace{1.8em}）}，最小是\mbox{（\hspace{1.8em}）}。\par
}
\NativeSlot{452}{36}{%
% p105_s04

\ExamQuestion{4.}{5}%
慢跑\(10\)千米大约消耗\(631\)千卡热量，那么慢跑\(1\)千米大约会消耗\mbox{（\hspace{1.8em}）}千卡热量。妈妈两周共跑了\(100\)千米，大约能消耗\mbox{（\hspace{1.8em}）}千卡热量。\par
}
\NativeSlot{502}{36}{%
% p105_s05

\ExamQuestion{5.}{5}%
一个等腰三角形样式的风筝，底角是\(42^\circ\)，它的顶角是\mbox{（\hspace{1.8em}）}°，如果按角分，这是个\mbox{（\hspace{1.8em}）}三角形。\par
}
\NativeSlot{552}{61}{%
% p105_s06

\ExamQuestion{6.}{5}%
有\(4\)根小棒（如右图），从中挑选\(3\)根，围成一个三角形，这个三角形的周长是\mbox{（\hspace{1.8em}）}厘米。\par
\BookFigure{figures/p105_s06-01.tex}

% p105_s07

\ExamQuestion{7.}{5}%
一个两位小数保留一位小数后约是\(3.8\)，这个小数最大是\mbox{（\hspace{1.8em}）}，最小是\mbox{（\hspace{1.8em}）}。\par
}
\NativeSlot{655.5}{54}{%
% p105_s08

\ExamQuestion{8.}{5}%
四年级一班有\(42\)名同学，其中会下象棋的有\(19\)人，会下围棋的有\(16\)人，两种棋都不会下的有\(10\)人。两种棋都会下的有\mbox{（\hspace{1.8em}）}人，只会下象棋不会下围棋的有\mbox{（\hspace{1.8em}）}人。\par
}
